\chapter{Progettazione di Algoritmi, Modelli Computazionali e Gerarchia di Memoria}
\label{cha:lecture-one}

\begin{center}
  \textbf{Lezione 1}\\
  \emph{Introduzione, modelli computazionali e gerarchia di memoria}\\[0.5em]
  Data: 15 settembre 2026
\end{center}

\section{Definizione Formale di un Algoritmo}
\label{sec:formal-algorithm}

Un \emph{algoritmo} è una sequenza finita e ordinata di istruzioni elementari non ambigue che, a partire da un insieme specificato di dati di input, produce un corrispondente insieme di dati di output in un lasso di tempo finito.

Nella formulazione classica presentata in \emph{The Art of Computer Programming}~\cite{knuth1997}, un algoritmo deve soddisfare cinque proprietà fondamentali:

\begin{description}
  \item[Finitudine.] L'algoritmo deve sempre terminare dopo un numero finito di passi di calcolo per ogni configurazione di input ammissibile.
  \item[Definitezza.] Ogni istruzione deve essere formulata in modo chiaro, preciso e senza ambiguità; lo stato di esecuzione successivo deve essere univocamente determinato.
  \item[Input.] L'algoritmo riceve zero o più quantità specificate dall'esterno prima che l'esecuzione abbia inizio.
  \item[Output.] L'algoritmo produce una o più quantità di output che hanno una relazione specificata con i dati di input. L'output deve essere significativo per il problema che si sta risolvendo.
  \item[Efficacia.] Ogni operazione prescritta deve essere elementare e fisicamente eseguibile in tempo finito utilizzando una quantità limitata di risorse, come carta e penna o una macchina sequenziale.
\end{description}

\subsection{Algoritmi Esatti, Intelligenza Artificiale ed Euristiche}
\label{sec:exact-vs-heuristic}

Un classico algoritmo esatto si distingue da un metodo di intelligenza artificiale o euristico per le garanzie che fornisce e per il modo in cui esplora lo spazio delle decisioni.

\begin{table}[htbp]
  \centering
  \caption{Confronto tra algoritmi esatti, intelligenza artificiale e metodi euristici.}
  \label{tab:algorithm-comparison}
  \begin{tabularx}{\textwidth}{@{}>{\raggedright\arraybackslash}p{0.20\textwidth}YY@{}}
    \toprule
    Dimensione & Algoritmo esatto classico & Metodo di IA o euristico \\
    \midrule
    Garanzia di correttezza & Produce una soluzione corretta per ogni istanza ammissibile; l'ottimalità è garantita quando la specifica del problema richiede un ottimo. & Può privilegiare la qualità empirica o l'efficienza pratica e può mancare di una garanzia universale per ogni istanza. \\
    Terminazione & Può essere dimostrata in anticipo utilizzando funzioni potenziale o varianti di ciclo. & Spesso legata a criteri di arresto euristici, come un numero massimo di epoche o la convergenza del gradiente. \\
    Spazio delle decisioni & Determinato da una procedura specificata in modo preciso, che può essere deterministica o randomizzata. & Esplorato tramite ricerca euristica, inferenza statistica o ottimizzazione di una funzione di perdita. \\
    \bottomrule
  \end{tabularx}
\end{table}

La \emph{programmazione euristica} non cerca una dimostrazione formale di ottimalità assoluta su ogni istanza ammissibile. Piuttosto, cerca un compromesso efficiente per problemi computazionalmente intrattabili, come i problemi NP-completi o NP-ardui.

\section{Misurazione dell'Efficienza e Analisi Asintotica}
\label{sec:asymptotic-analysis}

\subsection{Perché Non Contare il Numero Esatto di Istruzioni?}

Definire il costo esatto di esecuzione $T_A(n)$ come il numero effettivo di cicli di clock o di istruzioni elementari presenta diverse difficoltà:

\begin{enumerate}
  \item \textbf{Dipendenza dall'architettura:} il risultato dipenderebbe dalla specifica CPU, dall'architettura del set di istruzioni (ISA) e dalla frequenza di clock.
  \item \textbf{Dipendenza dalla toolchain:} varierebbe a seconda del compilatore, del livello di ottimizzazione e del sistema operativo.
  \item \textbf{Mancanza di generalità:} oscurerebbe la qualità intrinseca dell'idea algoritmica.
\end{enumerate}

Per questo motivo, l'analisi degli algoritmi utilizza la notazione asintotica. Essa astrae dai dettagli implementativi e dalle costanti moltiplicative, e studia il tasso di crescita della funzione di costo al tendere della dimensione dell'input all'infinito: $n \to \infty$.

\begin{description}
  \item[O-grande (limite superiore asintotico):]
    \[
      \begin{aligned}
        f(n) = \mathcal{O}(g(n)) &\iff
        \exists c > 0,\ n_0 > 0\text{ tali che}\\
        &0 \leq f(n) \leq c\,g(n),\quad \forall n \geq n_0.
      \end{aligned}
    \]
  \item[Omega-grande (limite inferiore asintotico):]
    \[
      \begin{aligned}
        f(n) = \Omega(g(n)) &\iff
        \exists c > 0,\ n_0 > 0\text{ tali che}\\
        &0 \leq c\,g(n) \leq f(n),\quad \forall n \geq n_0.
      \end{aligned}
    \]
  \item[Theta-grande (limite asintotico stretto):]
    \[
      f(n) = \Theta(g(n)) \iff
      f(n) = \mathcal{O}(g(n)) \land f(n) = \Omega(g(n)).
    \]
\end{description}

\subsection{Perché Preferire l'Analisi al Caso Pessimo?}

\begin{itemize}
  \item \textbf{Garanzia assoluta:} un limite superiore garantisce che l'algoritmo non impiegherà mai più tempo del limite teorico. Questo è essenziale per i sistemi real-time e mission-critical.
  \item \textbf{Semplice derivazione formale:} l'analisi al caso pessimo non richiede assunzioni probabilistiche sulla distribuzione dell'input, le quali sono spesso necessarie per l'analisi al caso medio.
\end{itemize}

\section{Modelli Computazionali: RAM e Hardware Reale}
\label{sec:ram-model}

Un algoritmo non viene eseguito in un vuoto astratto: viene eseguito su una macchina con processori, memorie e meccanismi di trasferimento dati. Un modello computazionale specifica quali dettagli della macchina sono rilevanti per l'analisi e quali vengono deliberatamente ignorati. Il modello RAM è il punto di partenza standard; l'hardware reale spiega poi quando tale astrazione diviene insufficiente.

\subsection{Il Modello Random Access Machine}

Il modello Random Access Machine (RAM) è un'astrazione sequenziale, basata su parole (word), ispirata all'architettura di von Neumann. Una RAM è costituita da un processore, un programma finito e una memoria indirizzabile di parole. Le sue assunzioni abituali sono:

\begin{itemize}
  \item \textbf{Esecuzione sequenziale:} un processore esegue un'istruzione elementare alla volta.
  \item \textbf{Dati della dimensione di una parola:} l'aritmetica, i confronti, le assegnazioni e il calcolo degli indirizzi su un numero costante di parole richiedono tempo costante.
  \item \textbf{Accesso uniforme alla memoria:} la lettura o la scrittura di qualsiasi indirizzo costa $\mathcal{O}(1)$, indipendentemente dall'indirizzo e dalla storia degli accessi.
  \item \textbf{Memoria sufficiente:} l'input e le strutture dati risiedono interamente nello spazio di indirizzamento modellato; il costo di allocazione dello spazio viene analizzato separatamente quando risulta rilevante.
\end{itemize}

Il modello RAM è utile perché rende confrontabili le idee algoritmiche. Ad esempio, un ciclo che esamina $n$ elementi di un array esegue $\Theta(n)$ operazioni RAM, indipendentemente dal particolare processore utilizzato per eseguirlo. Il modello non afferma che una macchina reale possieda memoria infinita o una latenza veramente costante; esso dichiara quali costi vengono astratti.

\subsection{Hardware Reale: Gerarchia di Memoria e Memory Wall}

Su una macchina moderna, il processore è molto più veloce delle memorie che contengono grandi insiemi di dati. Questo crescente divario tra calcolo e movimento dei dati è noto come \emph{memory wall}. Inoltre, la memoria locale non è una singola risorsa uniforme: è organizzata in livelli con capacità, latenze e granularità di trasferimento differenti.

Una gerarchia semplificata, ordinata dal livello più veloce e più piccolo a quello più lento e più grande, è:

\[
\begin{aligned}
  \text{Registri della CPU}
    &\longrightarrow \text{Cache L1}
    \longrightarrow \text{Cache L2}
    \longrightarrow \text{Cache L3}
    \longrightarrow \text{DRAM} \\
    &\longrightarrow \text{NVMe/SSD/disco}
    \longrightarrow \text{Storage remoto / cloud}.
\end{aligned}
\]

I registri e le cache fanno parte del sistema di memoria veloce del processore; la DRAM è l'usuale memoria principale; gli SSD e i dischi forniscono archiviazione persistente. Lo storage remoto e i servizi cloud sono inclusi qui solo come una più ampia gerarchia di accesso ai dati, in quanto sono raggiunti attraverso una rete piuttosto che attraverso il bus di memoria locale.

\begin{figure}[htbp]
  \centering
  \begin{tikzpicture}[
    level/.style={draw, rounded corners=2pt, minimum height=8mm,
      align=center, font=\small},
    link/.style={-{Latex[length=2.5mm]}, thick},
    node distance=6mm
  ]
    \node[level, fill=blue!18, minimum width=28mm] (registers)
      {Registri};
    \node[level, fill=blue!14, minimum width=43mm, below=of registers] (l1)
      {Cache L1};
    \node[level, fill=blue!11, minimum width=58mm, below=of l1] (l2)
      {Cache L2};
    \node[level, fill=blue!8, minimum width=73mm, below=of l2] (l3)
      {Cache L3};
    \node[level, fill=green!12, minimum width=88mm, below=of l3] (dram)
      {DRAM / memoria principale};
    \node[level, fill=orange!16, minimum width=103mm, below=of dram] (storage)
      {NVMe / SSD / disco};
    \node[level, fill=gray!18, minimum width=118mm, below=of storage] (network)
      {Rete / cloud};

    \draw[link] (registers) -- (l1);
    \draw[link] (l1) -- (l2);
    \draw[link] (l2) -- (l3);
    \draw[link] (l3) -- (dram);
    \draw[link] (dram) -- (storage);
    \draw[link] (storage) -- (network);

    \node[align=center, above=of registers, font=\footnotesize] (fast)
      {più veloce\\più piccolo\\più costoso per bit};
    \node[align=center, below=of network, font=\footnotesize] (large)
      {più lento\\più grande\\più economico per bit};
  \end{tikzpicture}
  \caption{Una vista gerarchica della memoria e dell'archiviazione.}
  \label{fig:memory-hierarchy}
\end{figure}

\autoref{fig:memory-hierarchy} descrive i principali livelli di memoria locale: ogni strato inferiore scambia una minore velocità di accesso per una maggiore capacità.

Spostandosi verso il basso nella \autoref{fig:memory-hierarchy}, la capacità generalmente cresce mentre crescono anche la latenza e il costo di trasferimento. I rapporti esatti dipendono dall'hardware e dal carico di lavoro, per cui valori come da $100$ a $200$ volte per la DRAM rispetto alla L1, o da $10^4$ a $10^6$ volte per il disco rispetto alla L1, dovrebbero essere letti come esempi di ordine di grandezza piuttosto che come costanti universali.

\subsection{Quando il Modello RAM Non È Sufficiente}

Il modello RAM è una buona prima approssimazione quando il calcolo è prevalentemente sequenziale, il working set risiede interamente nella memoria principale e il costo del movimento dei dati non è il fattore dominante. Un modello diverso è preferibile quando uno dei seguenti effetti controlla le prestazioni:

\begin{description}
  \item[Calcolo su GPU.] L'esecuzione SIMT (Single Instruction, Multiple Threads) utilizza molti core di elaborazione e distingue tra spazi di memoria condivisa, locale e globale.
  \item[Memoria esterna.] Nel calcolo out-of-core, il dataset non risiede interamente nella memoria veloce, per cui i trasferimenti a blocchi e le operazioni di I/O dominano il costo.
  \item[Calcolo parallelo e distribuito.] Diversi processori o macchine eseguono simultaneamente e la sincronizzazione e la comunicazione diventano costi espliciti.
  \item[Informatica quantistica.] Qubit, sovrapposizione, entanglement e porte unitarie richiedono un modello basato su operazioni quantistiche piuttosto che su parole RAM.
\end{description}

La lezione ingegneristica è quindi non di scartare il modello RAM, ma di utilizzare il modello più semplice che catturi ancora il collo di bottiglia del sistema target.

\section{Località e Modello a Due Livelli}
\label{sec:external-memory}

I sistemi moderni sfruttano due proprietà statistiche dei programmi reali:

\begin{enumerate}
  \item \textbf{Località temporale:} se una locazione di memoria viene referenziata al tempo $t$, è probabile che venga referenziata nuovamente a breve, al tempo $t + \Delta t$.
  \item \textbf{Località spaziale:} se una locazione di memoria $x$ viene referenziata, locazioni vicine come $x+1$, $x+2$ e così via, verranno probabilmente referenziate a breve.
\end{enumerate}

\subsection{Il Modello di Memoria a Due Livelli}

Quando si elaborano gigabyte o terabyte di dati, la gerarchia multilivello può essere astratta come un modello a due livelli, associato al modello a memoria esterna~\cite{aggarwal1988}, noto anche come modello Disk Access Machine.

\begin{figure}[htbp]
  \centering
  \begin{tikzpicture}[
    box/.style={draw, minimum width=38mm, minimum height=10mm, align=center},
    arrow/.style={-{Latex[length=2.5mm]}, thick},
    node distance=12mm
  ]
    \node[box] (cpu) {CPU};
    \node[box, below=of cpu] (fast) {Memoria veloce $M$\\capacità: $M$ parole};
    \node[box, below=of fast] (slow) {Memoria esterna $D$\\illimitata, blocchi di $B$ parole};
    \draw[arrow] (cpu) -- node[right, font=\small] {una parola, $\mathcal{O}(1)$} (fast);
    \draw[arrow] (fast) -- node[right, font=\small] {un blocco = $B$ parole} (slow);
  \end{tikzpicture}
  \caption{Modello di memoria a due livelli con memoria veloce $M$ e memoria esterna $D$.}
  \label{fig:two-level-memory}
\end{figure}

\autoref{fig:two-level-memory} descrive l'astrazione a due livelli: la CPU opera su parole nella memoria veloce $M$, mentre i trasferimenti tra $M$ e la memoria esterna $D$ avvengono in blocchi di $B$ parole.

\subsubsection{Parametri del Modello}

\begin{description}
  \item[$N$:] numero totale di elementi del problema da elaborare.
  \item[$M$:] dimensione della memoria veloce interna, misurata come il numero massimo di elementi che possono essere contenuti simultaneamente.
  \item[$B$:] dimensione del blocco trasferito in una singola operazione di I/O, con $M \geq 2B$.
\end{description}

\subsubsection{Metriche di Costo}

\begin{enumerate}
  \item \textbf{Complessità di I/O (metrica primaria):} numero di trasferimenti di blocchi di dimensione $B$ tra la memoria lenta e quella veloce.
  \item \textbf{Tempo di CPU (metrica secondaria):} numero di operazioni logiche e aritmetiche eseguite dopo che i dati sono stati caricati in $M$.
  \item \textbf{Spazio:} numero di blocchi allocati sulla memoria secondaria.
\end{enumerate}

Questa astrazione è frattale: può modellare sia i trasferimenti RAM--disco sia i trasferimenti cache L3--RAM, ridefinendo il significato di $M$ e $B$.

\section{Esempio Pratico: Scansione di un Array Lineare}
\label{sec:linear-scan}

Consideriamo un array contiguo di $N$ interi scansionato in sequenza, ad esempio per calcolare la somma $\sum_{i=1}^{N} A[i]$.

\begin{lstlisting}[style=code, language=C, caption={Scansione sequenziale di un array}, label={lst:linear-scan}]
long sum = 0;
for (int i = 0; i < N; i++) {
    sum += A[i];
}
\end{lstlisting}

\subsection{Analisi nel Classico Modello RAM}

\begin{itemize}
  \item Accessi in memoria: $N$.
  \item Istruzioni di addizione: $N$.
  \item Tempo di calcolo asintotico: $\Theta(N)$.
\end{itemize}

\subsection{Analisi nel Modello a Due Livelli}

Assumiamo che la memoria secondaria memorizzi l'array $A$ in modo consecutivo in blocchi di dimensione $B$ e che $A[0]$ sia allineato con il confine di un blocco. Senza questa assunzione di allineamento, la scansione potrebbe richiedere al più un trasferimento di blocco aggiuntivo, senza alterare il limite asintotico:

\begin{enumerate}
  \item Al primo accesso ad $A[0]$, si verifica un miss. Il sistema trasferisce il blocco contenente $A[0], \ldots, A[B-1]$ nella memoria veloce con un singolo I/O.
  \item I successivi $B-1$ accessi trovano i loro dati già nella memoria veloce, non richiedendo ulteriori I/O. Questo sfrutta la località spaziale.
  \item Un nuovo trasferimento di I/O è necessario solo quando l'indice attraversa il confine di un blocco, ovvero, una volta ogni $B$ elementi.
\end{enumerate}

Pertanto, il costo totale di I/O di una scansione lineare è
\[
  \operatorname{Scan}(N) = \left\lceil \frac{N}{B} \right\rceil
  = \Theta\!\left(\frac{N}{B}\right)\text{ operazioni di I/O}.
\]

\subsection{Accesso Lineare contro Accesso Sparso}

\begin{itemize}
  \item \textbf{Accesso sequenziale (alta località spaziale):} la lettura di $N$ elementi richiede $N/B$ operazioni di I/O. Poiché $B$ varia tipicamente da $64$ byte per una linea di cache a centinaia di kilobyte o megabyte per pagine e blocchi del disco, lo speedup rispetto al modello ingenuo può essere di un fattore approssimativamente pari a $B$.
  \item \textbf{Accesso casuale o passo $k>B$ (nessuna località spaziale):} ogni elemento richiesto potrebbe risiedere in un blocco diverso, causando $N$ operazioni di I/O indipendenti. Il costo degrada a $\Theta(N)$ I/O, il che è asintoticamente più lento di un fattore $B$.
\end{itemize}
